% !Mode:: "TeX:UTF-8"

\hitsetup{
  info={
    %******************************
    % 注意：
    %   1. 配置里面不要出现空行
    %   2. 不需要的配置信息可以删除
    %   3. 中英文两版的字段收在 zh={} 与 en={} 里，字段名不带后缀。
    %      这是推荐写法；带后缀的 title-zh={...} 也认，两种等价
    %******************************
    %
    %=====
    % 秘级
    %=====
    secret-level={公开},
    classified-index-national={TM301.2},
    classified-index-international={62-5},
    %
    %===========
    % 不分语言的
    %===========
    report-number={no9527}, %编号
    %
    %=========
    % 中文信息
    %=========
    zh={
      % 花括号分组给「题目栏是两条固定线」的封面分行（本科第二封面、博后第一封面）；
      % \\ 给自由排的封面分行；方括号是副标题，写了就排。
      % 带副标题写成 title={...}{...}[一条副标题]。副标题只排在自由排的封面上，
      % 题目栏是两条固定线的那几种（博后第一、威海本科第二、深圳博士中期）没有这一栏
      discipline={工学},
      associate-supervisor={某某某教授}, % 副指导老师
      co-supervisor={某某某教授}, % 联合指导老师
      % 日期不写就取编译当天。要指定就写 ISO 数字：yyyy、yyyy-mm、yyyy-mm-dd
      % 都认，月日补不补零都行；中文封面出「2026 年 6 月」，英文封面出
      % 「June, 2026」。写成别的样子就原样排出来，日志里会说一句。
      %date={2026-06},
      student-type={同等学力人员}, %非全日制教育申请学位者
      %（同等学力人员）、（工程硕士）、（工商管理硕士）、
      %（高级管理人员工商管理硕士）、（公共管理硕士）、（中职教师）、（高校教师）等
      station={哈铁西站}, %博士后站名称
      research-start-date={3050-09-10}, % 研究工作起始时间
      research-end-date={3090-10-10}, % 研究工作期满时间
      % 关键词用“英文逗号”分割
      keywords={\TeX, \LaTeX, CJK, 嗨！, thesis},
    },
    %
    %=========
    % 英文信息
    %=========
    en={
      % 带副标题写成 title={...}[This is the sub title]
      title={Research on key technologies of partial porous externally pressurized gas bearing},
      discipline={Engineering},
      major={Computer Science and Technology},
      affiliation={\emultiline[t]{School of Mechatronics Engineering \\ Mechatronics Engineering}},
      author={Yu Dongmei},
      supervisor={Professor XXX},
      associate-supervisor={XXX},
      % 不写就取编译当天。要指定就写 ISO 数字，模板按各处的语言和精度排：
      % 中文封面出「2017 年 12 月」，英文封面出「December, 2017」。
      % yyyy、yyyy-mm、yyyy-mm-dd 都认，月日补不补零都行。
      date={2017-12},
      student-type={Master of Art},
      keywords={\TeX, \LaTeX, CJK, template, thesis},
    },
  },
}

\begin{abstract-zh}

摘要的字数（以汉字计），硕士学位论文一般为500～1000字，博士学位论文为1000～2000字，
均以能将规定内容阐述清楚为原则，文字要精练，段落衔接要流畅。摘要页不需写出论文题目。
英文摘要与中文摘要的内容应完全一致，在语法、用词上应准确无误，语言简练通顺。
留学生的英文版博士学位论文中应有不少于3000字的“详细中文摘要”。

  关键词是为了文献标引工作、用以表示全文主要内容信息的单词或术语。关键词不超过 5
  个，每个关键词中间用分号分隔。（模板作者注：关键词分隔符不用考虑，模板会自动处
  理。英文关键词同理。）
\end{abstract-zh}

\begin{abstract-en}
   An abstract of a dissertation is a summary and extraction of research work
   and contributions. Included in an abstract should be description of research
   topic and research objective, brief introduction to methodology and research
   process, and summarization of conclusion and contributions of the
   research. An abstract should be characterized by independence and clarity and
   carry identical information with the dissertation. It should be such that the
   general idea and major contributions of the dissertation are conveyed without
   reading the dissertation.

   An abstract should be concise and to the point. It is a misunderstanding to
   make an abstract an outline of the dissertation and words ``the first
   chapter'', ``the second chapter'' and the like should be avoided in the
   abstract.

   Key words are terms used in a dissertation for indexing, reflecting core
   information of the dissertation. An abstract may contain a maximum of 5 key
   words, with semi-colons used in between to separate one another.
\end{abstract-en}
